%!TEX program = xelatex
\documentclass[10pt,a4paper,fontset=none]{ctexart}

\usepackage[a4paper,top=1.18cm,bottom=1.05cm,left=1.45cm,right=1.45cm]{geometry}
\usepackage{xcolor}
\usepackage{fontspec}
\usepackage{tabularx}
\usepackage{enumitem}
\usepackage{titlesec}
\usepackage{hyperref}
\usepackage{microtype}
\usepackage{ragged2e}
\usepackage{array}

\definecolor{msblue}{HTML}{0067B8}
\definecolor{ink}{HTML}{17212B}
\definecolor{muted}{HTML}{4E5968}
\definecolor{rulegray}{HTML}{D6DCE3}

\setmainfont{Avenir Next}
\setsansfont{Avenir Next}
\setCJKmainfont[
  Path=/usr/local/texlive/2023/texmf-dist/fonts/opentype/public/fandol/,
  BoldFont=FandolSong-Bold.otf
]{FandolSong-Regular.otf}
\setCJKsansfont[
  Path=/usr/local/texlive/2023/texmf-dist/fonts/opentype/public/fandol/,
  BoldFont=FandolHei-Bold.otf
]{FandolHei-Regular.otf}
\setCJKmonofont[
  Path=/usr/local/texlive/2023/texmf-dist/fonts/opentype/public/fandol/
]{FandolFang-Regular.otf}
\newCJKfontfamily\awardbrandfont{Heiti SC}
\renewcommand{\familydefault}{\sfdefault}
\color{ink}
% Preserve searchable/copyable Unicode text in the generated PDF.
\XeTeXgenerateactualtext=1

\hypersetup{
  colorlinks=true,
  urlcolor=msblue,
  linkcolor=msblue,
  pdfauthor={Yuwei Zhao},
  pdftitle={赵雨薇 - 中文简历}
}

\pagestyle{empty}
\setlength{\parindent}{0pt}
\setlength{\parskip}{0pt}
\setlist[itemize]{leftmargin=1.35em,itemsep=1.2pt,topsep=0pt,parsep=0pt,partopsep=0pt}

\titleformat{\section}
  {\large\bfseries\color{msblue}}
  {}{0pt}{}
  [\vspace{-3pt}\color{rulegray}\titlerule]
% Every module uses the same fixed space above and below its heading.
\titlespacing*{\section}{0pt}{6pt}{8pt}

\newcommand{\datedentry}[3]{%
  \begin{tabularx}{\textwidth}{@{}>{\bfseries}X>{\raggedleft\arraybackslash}p{3.2cm}@{}}
    #1 & \textcolor{muted}{#2}\\[-1pt]
    \multicolumn{2}{@{}l@{}}{\textcolor{muted}{#3}}
  \end{tabularx}
}

\newcommand{\paperentry}[3]{%
  \item \textbf{#1}\\[-1pt]
  {\small #2\hfill\textbf{\textcolor{msblue}{#3}}}
}

\newcommand{\researchpaperentry}[6]{%
  \item \textbf{#1}\\[-1pt]
  {\small #2\hfill\textbf{\textcolor{msblue}{#3}}}
  \begin{itemize}[leftmargin=1.3em,label={\textcolor{msblue}{\scriptsize\textbullet}},itemsep=0pt,topsep=1pt,parsep=0pt,partopsep=0pt]
    \small
    \item \textbf{方法：}#4
    \item \textbf{历史优化复现：}#5
    \item \textbf{真实项目优化：}#6
  \end{itemize}
}

\newcommand{\paperentrywithsummary}[4]{%
  \item \textbf{#1}\\[-1pt]
  {\small #2\hfill\textbf{\textcolor{msblue}{#3}}}
  \begin{itemize}[leftmargin=1.3em,label={\textcolor{msblue}{\scriptsize\textbullet}},itemsep=0pt,topsep=1pt,parsep=0pt,partopsep=0pt]
    \small
    \item \textbf{主要内容：}#4
  \end{itemize}
}

\newcommand{\awardentry}[2]{%
  \noindent{\awardbrandfont #1}\hfill\textcolor{muted}{#2}\par\vspace{2pt}
}

\begin{document}

\begin{tabularx}{\textwidth}{@{}X>{\raggedleft\arraybackslash}p{10.2cm}@{}}
  {\fontsize{25}{28}\selectfont\bfseries 赵雨薇}\quad
  {\large\color{muted}Yuwei Zhao}
  & {\small
      \begin{tabular}{@{}l@{\hspace{1.1em}}l@{}}
        \textbf{电话：}173 9804 7360 &
        \textbf{邮箱：}\href{mailto:zhaoyuwei@stu.pku.edu.cn}{zhaoyuwei@stu.pku.edu.cn}\\[-1pt]
        \multicolumn{2}{@{}l@{}}{\textbf{个人主页：}\href{https://fyrsta7.github.io/}{https://fyrsta7.github.io/}}
      \end{tabular}
    }\\
\end{tabularx}

\vspace{4pt}
{\color{msblue}\rule{\textwidth}{1.35pt}}\par

\section{教育经历}
\datedentry{北京大学｜计算机学院｜博士研究生}{2024.09 -- 至今}{导师：熊英飞教授}
\vspace{3pt}
\datedentry{北京大学｜信息科学技术学院｜本科}{2020.09 -- 2024.06}{专业：信息与计算科学｜毕业论文：基于抽象解释的动态规划算法自动合成}
\vspace{-10pt}

\section{科研概况}
研究方向为\textbf{基于大语言模型的自动代码性能优化}，探索如何利用大模型理解程序、识别优化机会并生成可验证的高性能实现。研究综合运用程序分析、程序综合、规则与样例引导、性能测试等技术，重点考察方法在多种编程语言和真实大型项目中的性能收益、可靠性与可迁移性。

\section{论文成果}
\begin{itemize}
  \researchpaperentry
    {Multi-Source and Cross-Scenario Strategy-Guided Code Optimization}
    {\textbf{Yuwei Zhao}, Qianyu Xiao, Ye Cui, Yijun Yu, Yingfei Xiong}
    {Under Review, 2026}
    {提出 MoST 多知识源、跨场景策略迁移框架，以证据聚类与迁移生成静态分析规则，突破单源知识与跨语言复用限制。}
    {在 151 个 C/C++、150 个 Python、50 个 Rust 历史优化任务中复现开发者修改，成功复现数较 SemOpt 提升 21.88\%--180.00\%。}
    {在 15 个真实项目中发现并应用原仓库尚未采用的新优化，单项最高、项目平均分别提升 19.72\%--717.42\%、4.44\%--258.17\%。相较 SemOpt 与 Codex，15/15 个项目单项提升最高，14/15 个项目平均提升最高。}
  \researchpaperentry
    {SemOpt: LLM-Driven Code Optimization via Rule-Based Analysis}
    {\textbf{Yuwei Zhao}, Yuan-An Xiao, Qianyu Xiao, Zhao Zhang, Yingfei Xiong}
    {TOSEM (CCF-A), 2026}
    {提出 LLM 与静态分析规则驱动的优化框架，从历史提交中提炼策略并生成 Semgrep 规则，克服语法检索难以匹配语义等价优化的局限。}
    {在 151 个 C/C++、150 个 Python 历史优化任务中复现开发者修改，成功复现数较基线分别提升 1.38--28 倍、4.60--6.33 倍。}
    {在 10 个大型项目中发现并应用原仓库尚未采用的新优化，单项最高、项目平均分别提升 5.04\%--479.90\%、1.68\%--143.25\%。}
  \paperentrywithsummary
    {Superfusion: Eliminating Intermediate Data Structures via Inductive Synthesis}
    {Ruyi Ji, \textbf{Yuwei Zhao}, Nadia Polikarpova, Yingfei Xiong, Zhenjiang Hu}
    {PLDI (CCF-A), 2024}
    {SuFu 以归纳程序综合消除函数式程序的中间数据结构，在 290 个任务中解决 264 个。}
  \paperentrywithsummary
    {Decomposition-Based Synthesis for Applying Divide-and-Conquer-Like Algorithmic Paradigms}
    {Ruyi Ji, \textbf{Yuwei Zhao}, Yingfei Xiong, Di Wang, Lu Zhang, Zhenjiang Hu}
    {TOPLAS (CCF-A), 2024}
    {AutoLifter 以组件消除和变量消除分解综合任务，自动应用 6 类分治算法范式，在 96 个任务中解决 82 个。}
\end{itemize}

\section{获奖经历}
\awardentry{2026年鲲鹏昇腾 KADC 明日之星先锋奖}{2026}
\awardentry{北京大学计算机学院九坤奖学金}{2025}
\awardentry{北京大学三好学生}{2025、2023、2021}
\awardentry{北京市普通高等学校优秀毕业生}{2024}
\awardentry{北京大学优秀毕业生}{2024}
\awardentry{北京大学三等奖学金}{2023、2021}
\awardentry{北京大学学习优秀奖}{2022}

\end{document}
